\chapter{Business Metrics Reference}\label{ch:metrics-catalog}

% Scope: structured pedagogical reference --- formula, term breakdown, what-it-tracks,
%   why-it-matters, benchmark, trap, and worked example for all 30 metrics.
%   Six sections: Unit Economics, SaaS & Recurring Revenue, Funnel & Marketing,
%   Ad-Platform, Product & Engagement, Vanity Metrics.
%   House style: flowing prose under headings; example and admonition set-offs.
%   No metriccard/mcfield/multicols in this file.

This chapter is a structured reference, not an argument. Each metric gets a definition, a
formula with every term explained, a plain-English statement of what the number actually
tracks, the downstream decision it informs, a defensible benchmark where one exists, and the
trap that makes the metric lie. Definitions follow \textcite{bendle2020marketing}; the
derivations live in the home chapters --- CAC and LTV in \cref{ch:customer-economics},
ROAS/CPA mechanics in \cref{ch:paid-acquisition-platforms}, recurring-revenue metrics in
\cref{ch:business-models}. A benchmark is an order-of-magnitude anchor, never a target in
itself.


%% ============================================================
\section{Unit Economics}

Unit economics answers one question before anything else: does a single transaction --- one
customer, one order, one contract --- make money, and by how much? Get the unit right and
scale is a lever; get it wrong and scale is an accelerant for losses.

\subsection{CAC --- Customer Acquisition Cost}

\textbf{Formula.}
\[
  \mathrm{CAC} = \frac{\text{Sales \& marketing spend}}{\text{New customers acquired}}
\]

\textbf{What each term means.}
\begin{description}
  \item[Sales \& marketing spend] All costs in a period that the business incurs to win new
    customers --- paid media, sales-team salaries (including benefits and on-costs),
    tooling subscriptions, content production, agencies, and attributable overhead. Fully
    loaded; \emph{not} only ad spend.
  \item[New customers acquired] Paying customers (not leads, trials, or signups) won in the
    same period used for the spend figure.
\end{description}

\textbf{What it tracks.} The fully loaded price the business pays to convert one prospect
into a paying customer.

\textbf{Why it matters.} CAC is one half of the acquisition-viability test. Compared against
LTV (\cref{eq:ltv-cac}), it caps how aggressively the firm can spend before unit economics
turn negative. A business that does not know its fully loaded CAC cannot know whether it is
creating or destroying value with every sale. The full derivation appears in
\cref{ch:customer-economics}.

\textbf{Benchmark.} Judge only against LTV (target LTV:CAC $\ge 3$). A standalone ``good
CAC'' does not exist --- the same \money{600} is excellent for an enterprise deal and ruinous
for a \money{10}/month consumer subscription.

\textbf{Trap.} Excluding sales and marketing salaries, tooling, and agency fees
understates CAC and silently inflates the LTV:CAC ratio --- the most common unit-economics
flattering trick in early-stage decks.

\begin{example}[CAC for the running SaaS]
Monthly spend in the running B2B SaaS: paid media \money{90000} + AE/SDR fully loaded
salaries \money{25000} + tooling \money{5000} = \money{120000}. New paying customers won in
the month: 200.
\[
  \mathrm{CAC} = \frac{\money{120000}}{200} = \money{600}.
\]
Strip out salaries and tooling and the reported figure falls to \money{450} --- a
\qty{25}{\percent} understatement that would flatter the 5.25$\times$ LTV:CAC ratio to
7$\times$, possibly triggering an acquisition-budget increase that destroys value at the true
unit economics.
\end{example}


\subsection{LTV --- Customer Lifetime Value}

\textbf{Formula.}
\[
  \mathrm{LTV} = \frac{m}{r - g}
\]
where this is the growing-perpetuity form (\cref{eq:gordon}) evaluated at the customer level.

\textbf{What each term means.}
\begin{description}
  \item[$m$] Annual contribution margin per customer --- revenue minus variable costs directly
    attributable to serving that customer. Use margin, not revenue; the formula measures
    economic value accruing to the firm, not topline.
  \item[$r$] The firm's discount rate (WACC or cost of equity) --- the opportunity cost of
    the capital tied up in serving customers. For the running SaaS, $r = \qty{11}{\percent}$.
  \item[$g$] The expected annual growth rate of that margin, capturing both price and volume
    expansion within the account. Must be strictly less than $r$ or the formula returns
    nonsense.
\end{description}

\textbf{What it tracks.} The present value of all future contribution margin a single
customer is expected to generate, discounted at the firm's cost of capital.

\textbf{Why it matters.} LTV sets the \emph{ceiling} on how much the firm should be willing
to pay to acquire a customer. It is mathematically identical to a discounted-cash-flow
terminal value (\cref{eq:gordon}) with the customer replacing the firm; the full derivation
is in \cref{ch:customer-economics}.

\textbf{Benchmark.} $\mathrm{LTV} \ge 3 \times \mathrm{CAC}$. This is not arbitrary: a 3$\times$
ratio leaves room for the cost of serving customers through their life (support, retention
spend, account management) without eroding the acquisition economics.

\textbf{Trap.} Plugging revenue rather than margin into $m$ overstates LTV dramatically ---
a \qty{42}{\percent} contribution-margin business that makes this error inflates LTV by
\qty{138}{\percent}. Equally, using a high growth rate $g$ that approaches $r$ can make LTV
arbitrarily large and render the metric useless for budgeting.

\begin{example}[LTV for the running SaaS]
ARPU is \money{600}/year; contribution margin is \qty{42}{\percent}, so $m =
\money{252}$/customer/year. Discount rate $r = \qty{11}{\percent}$; long-run margin growth
$g = \qty{3}{\percent}$.
\[
  \mathrm{LTV} = \frac{\money{252}}{0.11 - 0.03} = \frac{\money{252}}{0.08} = \money{3150}.
\]
If revenue had been plugged in instead of margin: $\money{600}/0.08 = \money{7500}$ --- a
\qty{138}{\percent} overstatement that would mislead the acquisition budget by the same
factor.
\end{example}


\subsection{LTV:CAC --- Lifetime Value to Acquisition Cost Ratio}

\textbf{Formula.}
\[
  \mathrm{LTV:CAC} = \frac{\mathrm{LTV}}{\mathrm{CAC}}
\]

\textbf{What each term means.}
\begin{description}
  \item[LTV] Customer lifetime value as defined above --- present value of future margin,
    not revenue.
  \item[CAC] Fully loaded customer acquisition cost as defined above.
\end{description}

\textbf{What it tracks.} Capital efficiency of acquisition: how many dollars of customer
value does the firm create for every dollar spent winning that customer?

\textbf{Why it matters.} LTV:CAC is the single most-cited unit-economics test in venture
and growth-equity underwriting. A ratio below 1 means the business loses money on every
customer before even counting fixed costs. The formula is fully derived as \cref{eq:ltv-cac}
in \cref{ch:customer-economics}.

\textbf{Benchmark.}
\begin{itemize}
  \item $\ge 3$: healthy; the standard institutional benchmark.
  \item 3--5: normal range for well-run SaaS.
  \item $> 5$: often signals under-investment in growth --- the firm is leaving value on
    the table by not acquiring customers it could profitably serve.
  \item $< 1$: each new customer destroys value; fix unit economics before scaling.
\end{itemize}

\textbf{Trap.} An LTV:CAC above 5 is not simply ``very good'' --- it typically means the
firm is spending too little on acquisition relative to the opportunity. Treating a high ratio
as a reward rather than a signal to increase spending is a common growth-phase error.

\begin{example}[LTV:CAC for the running SaaS]
From the calculations above, LTV $= \money{3150}$ and CAC $= \money{600}$.
\[
  \mathrm{LTV:CAC} = \frac{\money{3150}}{\money{600}} = 5.25.
\]
This sits above the 3--5 healthy range, suggesting the firm should test increasing acquisition
spend --- provided the incremental customer has the same LTV profile as the average.
\end{example}


\subsection{CAC Payback --- Months to Payback}

\textbf{Formula.}
\[
  \text{CAC Payback} = \frac{\mathrm{CAC}}{\text{Monthly margin per customer}}
\]

\textbf{What each term means.}
\begin{description}
  \item[CAC] Fully loaded acquisition cost per customer (\money{600} in the running SaaS).
  \item[Monthly margin per customer] Monthly contribution margin --- monthly ARPU multiplied
    by the contribution margin rate. This is the cash the customer actually throws off each
    month to recoup the acquisition investment.
\end{description}

\textbf{What it tracks.} How many months must a customer remain before the cumulative margin
from that customer repays the cost of acquiring them. Shorter payback periods reduce the
firm's cash-cycle risk and funding requirements.

\textbf{Why it matters.} Payback period translates LTV:CAC into a cash-flow reality check.
A high LTV:CAC is comforting but meaningless if customers churn before the payback month.
It is also the primary driver of how much capital the firm needs to sustain a given growth
rate; a business growing at \qty{20}{\percent} per month with a 24-month payback is far more
cash-intensive than one with an 8-month payback. The full derivation is \cref{eq:cac-payback}
in \cref{ch:customer-economics}.

\textbf{Benchmark.}
\begin{itemize}
  \item $< 12$ months: healthy for SMB-focused SaaS.
  \item $< 18$ months: acceptable for mid-market.
  \item $< 24$ months: borderline for enterprise; above 24 months demands strong retention
    data to justify the cash cycle.
\end{itemize}

\textbf{Trap.} The formula counts every customer in the denominator, including those who
churn before payback month. If monthly churn is high, the average customer never reaches
payback, making the stated figure an optimistic fiction. Cross-check payback against the
inverse of monthly churn (the average customer lifespan in months).

\begin{example}[CAC Payback for the running SaaS]
Monthly ARPU $= \money{50}$; contribution margin $= \qty{42}{\percent}$; monthly margin per
customer $= \money{50} \times 0.42 = \money{21}$. CAC $= \money{600}$.
\[
  \text{CAC Payback} = \frac{\money{600}}{\money{21}} \approx 29 \text{ months}.
\]
At monthly churn of \qty{0.65}{\percent}, average customer lifespan is approximately
$1/0.0065 \approx 154$ months, so almost all customers survive to payback. However, the 29-month
figure itself is above the 24-month enterprise threshold --- a signal to either reduce CAC
or find margin expansion levers (upsell, price increase) to compress payback.
\end{example}


\subsection{Contribution Margin}

\textbf{Formula.}
\[
  \text{Contribution Margin} = \text{Price} - \text{Variable cost per unit}
\]
Expressed as a rate: $\text{CM\%} = (\text{Price} - \text{Variable cost})/\text{Price}$.

\textbf{What each term means.}
\begin{description}
  \item[Price] Revenue per unit (per subscription seat, per transaction, per item sold).
  \item[Variable cost per unit] Costs that scale directly with each additional unit served:
    payment processing, cloud hosting increments, per-seat third-party APIs, direct material
    and labour for physical goods. Fixed costs (rent, base salaries, annual licences) are
    \emph{excluded}.
\end{description}

\textbf{What it tracks.} The cash each additional unit throws off after variable costs are
covered --- the pool of money available to cover fixed costs and ultimately generate profit.

\textbf{Why it matters.} Contribution margin is the economic engine of the unit. A
positive contribution margin means each incremental sale moves the firm toward breakeven; a
negative one means selling more accelerates losses regardless of volume. It feeds directly
into LTV (the $m$ term) and into CAC payback calculations. The relationship to gross margin
and operating profit is developed in \cref{ch:business-models}.

\textbf{Benchmark.} Positive and rising with scale as fixed-cost leverage improves. For
software, \qty{60}{\percent}+ is typical at the unit level before platform overhead; for
physical goods, margins below \qty{30}{\percent} make the unit-economics math very tight.

\textbf{Trap.} Confusing contribution margin with gross margin. Gross margin deducts cost
of goods sold as reported on the income statement (which may include semi-fixed items like
depreciation of production equipment); contribution margin is a management-accounting concept
that isolates only the truly variable costs. They converge for pure-software businesses but
diverge materially for any product with physical components or large fixed-production
infrastructure.

\begin{example}[Contribution Margin in the running SaaS]
Monthly subscription price: \money{50}. Variable costs per seat: payment processing
\money{1.50} + incremental cloud hosting \money{1.75} + per-seat API costs \money{0.75}
= \money{4.00} total.
\[
  \text{CM} = \money{50} - \money{4} = \money{46}; \quad
  \text{CM\%} = \frac{\money{46}}{\money{50}} = \qty{92}{\percent}.
\]
Wait --- why does the fact-sheet lock use \qty{42}{\percent}? Because the running SaaS
attributes a share of customer-success salaries and onboarding tooling as variable costs in
its management accounts (a legitimate fully-loaded treatment), raising variable cost per seat
to \money{29}. The \qty{42}{\percent} figure (\money{21}/month) is the number that flows into
LTV and CAC payback throughout this book.
\end{example}


\subsection{Gross Margin}

\textbf{Formula.}
\[
  \text{Gross Margin} = \frac{\text{Revenue} - \text{COGS}}{\text{Revenue}}
\]

\textbf{What each term means.}
\begin{description}
  \item[Revenue] Total recognised revenue in the period.
  \item[COGS (Cost of Goods Sold)] The direct costs of delivering the product or service as
    reported under GAAP/IFRS: direct labour, direct materials, hosting, third-party licences
    embedded in the product, amortisation of capitalised development costs directly tied to
    revenue delivery.
  \item[Gross Margin (as a ratio)] The share of each revenue dollar left after direct delivery
    costs; the starting point for all downstream profitability analysis.
\end{description}

\textbf{What it tracks.} The structural profitability of the firm's core offering before
operating expenses --- sales, marketing, R\&D, and G\&A. It answers: does the product itself
generate enough surplus to sustain the business's operating structure?

\textbf{Why it matters.} Gross margin is the structural ceiling on operating profitability.
A firm with \qty{40}{\percent} gross margin cannot, by arithmetic, reach a \qty{45}{\percent}
EBITDA margin --- no amount of operating efficiency overcomes that ceiling. For SaaS
investors, gross margin signals scalability: high margins ($>\qty{70}{\percent}$) indicate the
business can grow revenue without proportionally growing delivery costs.

\textbf{Benchmark.}
\begin{itemize}
  \item $>\qty{70}{\percent}$: typical healthy range for pure software.
  \item $>\qty{80}{\percent}$: elite; characteristic of high-value enterprise SaaS.
  \item \qty{40}--\qty{60}{\percent}: common for software with significant managed-service or
    implementation components.
  \item $<\qty{30}{\percent}$: typical for physical-goods or marketplace businesses where
    unit-level leverage is lower.
\end{itemize}

\textbf{Trap.} Burying hosting costs, support-team salaries, or third-party API fees in
operating expenses (rather than COGS) flatters gross margin on the income statement but
misleads anyone modelling scalability. This reclassification is a standard pre-fundraise
preparation in some startups; a sceptical analyst should reconstruct COGS from the cost
schedule, not accept the headline figure.

\begin{example}[Gross Margin in the running SaaS]
ARR $= \money{4000000}$ (MRR \money{340000} $\times$ 12 approximately). COGS as reported:
hosting \money{480000} + third-party APIs \money{240000} + amortisation of capitalised
development \money{120000} = \money{840000}.
\[
  \text{Gross Margin} = \frac{\money{4000000} - \money{840000}}{\money{4000000}}
    = \frac{\money{3160000}}{\money{4000000}} = \qty{79}{\percent}.
\]
If support-team salaries (\money{400000}) were also properly classified in COGS, gross margin
would fall to $(\money{4000000} - \money{1240000})/\money{4000000} = \qty{69}{\percent}$ ---
still healthy but below the elite threshold, changing the investor story.
\end{example}


%% ============================================================
\section{SaaS \& Recurring Revenue}

Recurring-revenue metrics measure the health of the subscription engine. They are
leading indicators: they reflect what the customer base is doing today and predict what the
income statement will show in twelve months. A founder who watches only revenue without these
metrics is reading a rearview mirror.

\subsection{MRR --- Monthly Recurring Revenue}

\textbf{Formula.}
\[
  \mathrm{MRR} = \sum_{\text{active subscriptions}} \text{monthly subscription value}
\]

\textbf{What each term means.}
\begin{description}
  \item[Monthly subscription value] The normalised recurring charge per subscription per
    month. Annual contracts are divided by 12; multi-year contracts by their total months.
    Usage fees, one-off professional-services revenue, and hardware sales are excluded.
  \item[$\sum$ over active subscriptions] Every subscription that is active (past the trial,
    not churned, not paused) at the measurement date.
\end{description}

\textbf{What it tracks.} The business's predictable monthly run-rate from subscription
contracts. It is the pulse of a recurring-revenue business, updated monthly as new customers
are added, existing customers expand or contract, and churned customers exit.

\textbf{Why it matters.} MRR is the foundation of all SaaS financial modelling: ARR,
NRR, burn multiple, and the Rule of 40 all derive from it. Investors use MRR growth rate ---
not revenue growth --- as the primary growth signal for pre-profit SaaS companies. The full
MRR waterfall model (new + expansion $-$ contraction $-$ churn) is derived as
\cref{eq:mrr-growth} in \cref{ch:business-models}.

\textbf{Benchmark.} The growth rate matters more than the level. A \money{340000} MRR
growing at \qty{15}{\percent}/month is fundamentally different from a \money{5000000} MRR
growing at \qty{1}{\percent}/month, even though the latter is larger.

\textbf{Trap.} Counting non-recurring revenue --- one-off implementation fees, variable
usage overages, or professional-services engagements --- as MRR. This inflates the metric
and, worse, creates false smoothness: when those one-off projects end, reported MRR drops and
the business looks like it is churning.

\begin{example}[MRR decomposition for the running SaaS]
The firm has \num{6700} paying customers at \money{50}/month average.
\[
  \mathrm{MRR} = \num{6700} \times \money{50} = \money{335000}.
\]
The fact-sheet lock uses \money{340000} MRR, reflecting the actual customer count of
\num{6800} at the measurement date (growth in the cohort from the previous month's
\num{6700}). One-off implementation fees of \money{45000} this month are excluded; including
them would overstate MRR by \qty{13}{\percent}.
\end{example}


\subsection{ARR --- Annual Recurring Revenue}

\textbf{Formula.}
\[
  \mathrm{ARR} = \mathrm{MRR} \times 12
\]

\textbf{What each term means.}
\begin{description}
  \item[MRR] Monthly recurring revenue as defined above --- subscriptions only.
  \item[$\times 12$] Annualisation assumes the current run-rate holds for twelve months.
    It is a \emph{snapshot}, not a forecast.
\end{description}

\textbf{What it tracks.} The annualised recurring run-rate of the business at a point in
time. ARR is the standard unit for SaaS valuation multiples (ARR multiple, derived as
\cref{eq:arr-multiple} in \cref{ch:business-models}) and for year-over-year growth
benchmarking.

\textbf{Why it matters.} ARR converts the MRR snapshot into the language investors and
acquirers use for comparables. At the \money{4000000} ARR level, the running SaaS is in the
range where institutional growth investors typically begin serious diligence; below
\money{1000000} ARR, most Series A institutional funds consider the data set too small to
extrapolate.

\textbf{Trap.} Annualising non-recurring revenue. If a business does \money{500000} in
one-off consulting this month and reports ARR of \money{6000000} by multiplying total revenue
by 12, the ``ARR'' will disappear next month. Investors who discover the error revalue the
business to its true recurring base.

\begin{example}[ARR for the running SaaS]
MRR $= \money{340000}$.
\[
  \mathrm{ARR} = \money{340000} \times 12 = \money{4080000}.
\]
The fact-sheet lock uses \money{4000000} ARR, representing a \money{333333} MRR baseline ---
the minor rounding reflects the normalised figure used for modelling. If the firm had
included last month's implementation fees, reported ARR would have been \money{4620000} --- a
\qty{15}{\percent} overstatement with no recurring substance.
\end{example}


\subsection{Logo Churn --- Customer Churn Rate}

\textbf{Formula.}
\[
  \text{Logo Churn} = \frac{\text{Customers lost in period}}{\text{Customers at start of period}}
\]

\textbf{What each term means.}
\begin{description}
  \item[Customers lost in period] Subscriptions that cancelled or did not renew during the
    measurement period (typically monthly or annual). Excludes downgrades (which are
    contraction, not churn).
  \item[Customers at start of period] The cohort base. Measure at the opening of the period,
    before new additions.
\end{description}

\textbf{What it tracks.} The rate at which the business loses customers by count, regardless
of their revenue size.

\textbf{Why it matters.} Logo churn is the most intuitive attrition signal: if
\qty{5}{\percent} of customers leave each month, the entire base turns over in roughly
20 months. At \qty{0.65}{\percent} monthly churn, the running SaaS's average customer
lifespan is $1/0.0065 \approx 154$ months --- over 12 years --- a very strong retention
story. Churn feeds directly into LTV: every percentage-point improvement in monthly churn
extends customer lifespan and raises LTV non-linearly through the perpetuity denominator.
The retention derivation is \cref{eq:retention} in \cref{ch:customer-economics}.

\textbf{Benchmark.}
\begin{itemize}
  \item $<\qty{1}{\percent}$/month (\qty{12}{\percent}/year): acceptable for SMB-focused
    SaaS.
  \item $<\qty{0.5}{\percent}$/month (\qty{6}{\percent}/year): good for mid-market.
  \item $<\qty{0.2}{\percent}$/month: elite enterprise retention.
  \item Annual contracts naturally suppress monthly churn figures; compare on a consistent
    basis.
\end{itemize}

\textbf{Trap.} Ignoring that lost logos may be tiny accounts. If the ten largest customers
(each at \money{10000}/month) leave but two hundred small accounts (each at \money{50}/month)
are added, logo churn looks healthy while revenue churn is catastrophic. Always pair logo
churn with net revenue retention (\cref{sec:saas-mrr} below) to get the full picture.

\begin{example}[Logo Churn in the running SaaS]
Opening customer count: \num{6700}. Cancellations in the month: 44.
\[
  \text{Logo Churn} = \frac{44}{\num{6700}} = \qty{0.657}{\percent}/\text{month}.
\]
The fact-sheet lock uses \qty{0.65}{\percent}/month. At this rate, average customer life
$\approx 154$ months; CAC payback of 29 months is well within that window.
\end{example}


\subsection{NRR --- Net Revenue Retention}

\textbf{Formula.}
\[
  \mathrm{NRR} = \frac{\text{MRR}_\text{start} + \text{Expansion MRR} - \text{Contraction MRR} - \text{Churned MRR}}{\text{MRR}_\text{start}}
\]

\textbf{What each term means.}
\begin{description}
  \item[MRR$_\text{start}$] Monthly recurring revenue from the existing customer base at the
    start of the period --- no new logos included.
  \item[Expansion MRR] Additional MRR from existing customers via upsell, cross-sell, or
    seat expansion.
  \item[Contraction MRR] MRR lost from existing customers who downgraded their plan without
    cancelling.
  \item[Churned MRR] MRR from customers who cancelled entirely.
\end{description}

\textbf{What it tracks.} The percentage of starting revenue the business retains from its
existing customer base after twelve months, net of expansions, contractions, and
cancellations.

\textbf{Why it matters.} NRR above \qty{100}{\percent} means the existing customer base
grows revenue even with zero new customer acquisition. This is the defining characteristic
of a compounding SaaS business: growth does not require perpetually refilling a leaky bucket.
At $\mathrm{NRR} = \qty{120}{\percent}$, a firm starting with \money{4000000} ARR would
reach \money{4800000} ARR from its current base alone in twelve months. See
\cref{ch:business-models} for the full mechanics.

\textbf{Benchmark.}
\begin{itemize}
  \item $>\qty{100}{\percent}$: net expansion --- the existing base grows itself.
  \item $>\qty{110}{\percent}$: good; mid-market SaaS benchmark.
  \item $>\qty{120}{\percent}$: elite; characteristic of the best enterprise SaaS companies.
  \item $<\qty{100}{\percent}$: the base is shrinking even before logo churn; a serious
    red flag requiring immediate attention.
\end{itemize}

\textbf{Trap.} A high NRR can mask heavy logo churn. If the business is losing many small
customers but the few large customers are expanding, NRR looks strong while the customer base
is concentrating into a small number of large accounts --- a fragility that is invisible in
the NRR figure alone.

\begin{example}[NRR for the running SaaS]
Starting MRR from existing \num{6700} customers: \money{335000}. Expansion (upsell to higher
tier): \money{12000}. Contraction (downgrades): \money{3500}. Churned MRR (44 customers $\times$
\money{50}): \money{2200}.
\[
  \mathrm{NRR} = \frac{\money{335000} + \money{12000} - \money{3500} - \money{2200}}
    {\money{335000}}
  = \frac{\money{341300}}{\money{335000}} \approx \qty{102}{\percent}.
\]
The business is net-expanding --- but modestly. Doubling the expansion motion (pushing NRR
to \qty{104}{\percent}) would add roughly \money{13000}/month in organic ARR growth from
the existing base alone.
\end{example}


\subsection{Rule of 40}

\textbf{Formula.}
\[
  \text{Rule of 40} = \text{Revenue growth rate (\%)} + \text{Profit margin (\%)}
\]

\textbf{What each term means.}
\begin{description}
  \item[Revenue growth rate] Year-over-year ARR or revenue growth expressed as a percentage.
  \item[Profit margin] Typically EBITDA margin or, for earlier-stage companies, free-cash-flow
    margin. The two are not interchangeable; clarify which is used before comparing across
    companies.
\end{description}

\textbf{What it tracks.} The balance between growth and profitability. A company spending
aggressively to grow (negative profit margin) can still satisfy the Rule of 40 if its growth
rate is high enough, and vice versa.

\textbf{Why it matters.} The Rule of 40 is the primary heuristic SaaS investors use to
evaluate the growth-efficiency trade-off at scale. It acknowledges that the ``right'' mix of
growth and margin is not fixed --- a \qty{60}{\percent}-growth company at $-\qty{20}{\percent}$
margins and a \qty{10}{\percent}-growth company at \qty{30}{\percent} margins score equally,
and both are acceptable in different phases. The formal derivation is \cref{eq:rule40} in
\cref{ch:business-models}.

\textbf{Benchmark.} $\ge 40$: the threshold at which investors typically assign premium
multiples. Above 60 is elite. Below 40 raises questions about the sustainability of the
growth-margin profile.

\textbf{Trap.} Gaming the score with one-time margin items (a large asset sale, a
restructuring credit, a deferred-revenue catch-up) that inflate profit margin in a single
period without changing the underlying business. Always examine the trailing twelve months
across multiple quarters to smooth one-offs.

\begin{example}[Rule of 40 for the running SaaS]
ARR grew from \money{3000000} to \money{4000000} over twelve months: growth rate
$= \qty{33}{\percent}$. EBITDA margin for the year: \qty{9}{\percent} (NOPAT-adjacent,
consistent with \money{600000} NOPAT on \money{4000000} revenue, \qty{15}{\percent}
NOPAT margin approximated at \qty{9}{\percent} EBITDA after adjustments).
\[
  \text{Rule of 40} = 33 + 9 = 42.
\]
The business clears the threshold --- barely. To reach 50 without accelerating growth, the
business would need to expand EBITDA margin to \qty{17}{\percent}, which at current revenue
means \money{680000} in additional operating profit --- achievable only through pricing power
or fixed-cost leverage, not incremental volume at current margins.
\end{example}


\subsection{Burn Multiple}

\textbf{Formula.}
\[
  \text{Burn Multiple} = \frac{\text{Net cash burn}}{\text{Net new ARR}}
\]

\textbf{What each term means.}
\begin{description}
  \item[Net cash burn] Cash consumed by the business in the period --- operating cash outflow
    net of any cash inflows from operations (collections, not fundraising).
  \item[Net new ARR] The incremental ARR added in the same period: new logo ARR + expansion
    ARR $-$ churned ARR $-$ contraction ARR. This is the ARR the burn ``bought.''
\end{description}

\textbf{What it tracks.} How many dollars of cash the business burns to generate each dollar
of net new ARR. It is a capital efficiency metric for growth --- lower is better.

\textbf{Why it matters.} Burn multiple connects the income statement to the balance sheet:
a business can be growing ARR rapidly while burning cash at a rate that will exhaust the
runway before the next funding event. At \$1 burn per \$1 new ARR, a firm is paying one
dollar of capital for one dollar of forward revenue --- efficient. At \$5 per \$1, it is
spending \$5 to acquire \$1 of ARR, which implies either a very low-quality growth motion
or a structural unit-economics problem. The formal definition is \cref{eq:burn-multiple} in
\cref{ch:business-models}.

\textbf{Benchmark.}
\begin{itemize}
  \item $< 1$: excellent capital efficiency; the business is near cash-flow positive from
    a growth perspective.
  \item $< 1.5$: good; most well-run growth-stage SaaS.
  \item $1.5$--$3$: acceptable in a deliberate land-grab phase with strong NRR.
  \item $> 3$: concerning; suggests growth is expensive relative to its quality.
\end{itemize}

\textbf{Trap.} The metric reads worse in deliberate land-grab phases where the firm is
investing ahead of revenue in market coverage, engineering, or customer success headcount.
Context matters: a \$2.5 burn multiple alongside \qty{150}{\percent} NRR is a fundamentally
different situation from a \$2.5 multiple alongside \qty{90}{\percent} NRR.

\begin{example}[Burn Multiple for the running SaaS]
Net cash burn in Q1: \money{450000}. Net new ARR in Q1: new logos (\num{200} customers
$\times$ \money{600}/year) = \money{120000}; expansion \money{48000}; churn loss \money{26400};
net new ARR $= \money{141600}$.
\[
  \text{Burn Multiple} = \frac{\money{450000}}{\money{141600}} \approx 3.18.
\]
This sits above the ``good'' threshold. Decomposing the burn reveals \money{120000} in R\&D
headcount hired ahead of an enterprise feature release --- if that release lifts NRR by
4 percentage points, the effective forward multiple will fall to approximately 1.8.
\end{example}


%% ============================================================
\section{Funnel \& Marketing}

Funnel metrics track the conversion of attention into revenue at each stage of the customer
journey. No single funnel metric should be optimised in isolation --- a conversion rate that
rises while average order value falls may represent a net loss; a falling CPA that brings in
lower-quality leads erodes LTV silently.

\subsection{CTR --- Click-Through Rate}

\textbf{Formula.}
\[
  \mathrm{CTR} = \frac{\text{Clicks}}{\text{Impressions}}
\]

\textbf{What each term means.}
\begin{description}
  \item[Clicks] The number of times a user clicked on the ad or link in the measurement
    period.
  \item[Impressions] The number of times the ad or link was shown to users (including to
    users who did not click).
\end{description}

\textbf{What it tracks.} The proportion of people who saw the ad or content and chose to
engage. CTR is a direct signal of creative and targeting resonance --- whether the message
is relevant enough to interrupt attention.

\textbf{Why it matters.} CTR is an early-funnel diagnostic. High CTR indicates the
creative, headline, or offer is compelling \emph{to the audience being shown it}. It does
not indicate quality of downstream conversion. Monitoring CTR alongside conversion rate
(below) reveals whether a campaign's problem is at the attention stage or the landing/offer
stage. Paid-platform CTR also feeds into Quality Score and eCPC mechanics, developed in
\cref{ch:paid-acquisition-platforms}.

\textbf{Benchmark.} Varies widely by channel and format; search ads typically run
\qty{2}{\percent}--\qty{6}{\percent}, display ads \qty{0.1}{\percent}--\qty{0.3}{\percent},
email links \qty{1}{\percent}--\qty{3}{\percent}. Compare only within channels and against
your own historical baseline.

\textbf{Trap.} High CTR with low downstream conversion rate means the ad is attracting the
wrong audience --- or promising something the landing page does not deliver. Optimising for
CTR alone (by writing misleading or sensationalist headlines) is a reliable way to destroy
CAC efficiency while hitting a vanity metric.

\begin{example}[CTR in the running SaaS]
A LinkedIn campaign targeting operations directors: \num{240000} impressions;
\num{4800} clicks.
\[
  \mathrm{CTR} = \frac{4800}{240000} = \qty{2}{\percent}.
\]
Of those \num{4800} clicks, \num{96} became trial signups (conversion rate \qty{2}{\percent})
and \num{20} became paying customers. The CTR looks healthy but the click-to-trial rate
reveals a landing-page problem --- the offer is resonating in the feed but not converting
once reached.
\end{example}


\subsection{CPC --- Cost Per Click}

\textbf{Formula.}
\[
  \mathrm{CPC} = \frac{\text{Total ad spend}}{\text{Total clicks}}
\]

\textbf{What each term means.}
\begin{description}
  \item[Total ad spend] The amount paid to the platform for the campaign in the measurement
    period.
  \item[Total clicks] The number of clicks delivered by the platform in the same period.
\end{description}

\textbf{What it tracks.} The effective price the business pays for a single visit from a
paid ad. CPC is set by the platform's auction and the advertiser's Quality Score (the
\cref{eq:ecpc} mechanism from \cref{ch:paid-acquisition-platforms}).

\textbf{Why it matters.} CPC is the input cost for all paid-traffic economics. At a given
conversion rate and average order value, CPC determines whether paid traffic is profitable.
The relationship is: $\mathrm{CPA} = \mathrm{CPC} / \mathrm{CVR}$, so a high CPC requires
either a high conversion rate or a high CLV to justify.

\textbf{Benchmark.} No universal benchmark; compare against the CPA ceiling implied by LTV.
Google Search in B2B SaaS typically runs \money{15}--\money{60} per click for competitive
terms; LinkedIn can run \money{8}--\money{25} per click.

\textbf{Trap.} Optimising for the lowest possible CPC rather than the lowest CPA. CPC can
be reduced by bidding on cheaper, lower-intent keywords --- but if those clicks convert at
one-quarter the rate, CPA rises even as CPC falls. Always evaluate CPC in the context of
downstream conversion.

\begin{example}[CPC in the running SaaS]
LinkedIn campaign: spend \money{12000}; \num{4800} clicks.
\[
  \mathrm{CPC} = \frac{\money{12000}}{4800} = \money{2.50}.
\]
With a click-to-paid conversion rate of $20/4800 = \qty{0.42}{\percent}$:
$\mathrm{CPA} = \money{2.50}/0.0042 \approx \money{595}$ --- almost exactly the \money{600}
CAC. This campaign is running at break-even on unit economics.
\end{example}


\subsection{CPM --- Cost Per Thousand Impressions}

\textbf{Formula.}
\[
  \mathrm{CPM} = \frac{\text{Total ad spend}}{\text{Total impressions}} \times 1000
\]

\textbf{What each term means.}
\begin{description}
  \item[Total ad spend] Amount paid to the platform in the measurement period.
  \item[Total impressions] Number of times the ad was shown (served) to users.
  \item[$\times 1000$] Normalises the figure to a per-thousand basis --- the standard
    industry unit for comparing reach costs across channels.
\end{description}

\textbf{What it tracks.} The cost of reaching one thousand viewers with an ad impression.
CPM is the primary buying and pricing unit for brand/awareness campaigns where
reach and frequency matter more than direct conversion.

\textbf{Why it matters.} CPM governs the cost of building top-of-funnel awareness. When
a brand runs TV, display, streaming audio, or programmatic display, CPM is the metric that
determines how much audience reach a given budget buys. In performance marketing, CPM is a
derivative of CTR and CPC rather than a primary optimisation target.

\textbf{Benchmark.} Programmatic display: \money{2}--\money{8} CPM. Social feeds:
\money{6}--\money{20} CPM. LinkedIn: \money{30}--\money{80} CPM (premium B2B targeting
commands a large premium). Connected TV: \money{20}--\money{50} CPM.

\textbf{Trap.} Optimising for cheap CPM on a worthless audience. A \money{1} CPM campaign
that reaches users with no interest in the product wastes every impression; a \money{40} CPM
LinkedIn campaign reaching the exact decision-maker persona may be dramatically more
efficient on a cost-per-outcome basis.

\begin{example}[CPM in the running SaaS]
A programmatic display retargeting campaign: spend \money{3000}; impressions \num{500000}.
\[
  \mathrm{CPM} = \frac{\money{3000}}{500000} \times 1000 = \money{6}.
\]
Clicks: \num{750} (CTR \qty{0.15}{\percent}); paid conversions: 3.
$\mathrm{CPA} = \money{3000}/3 = \money{1000}$ --- significantly above CAC. The cheap CPM
is purchasing worthless impressions; the retargeting audience needs tightening.
\end{example}


\subsection{CPA / CPL --- Cost Per Acquisition / Cost Per Lead}

\textbf{Formula.}
\[
  \mathrm{CPA} = \frac{\text{Ad spend}}{\text{Conversions}}
\]

\textbf{What each term means.}
\begin{description}
  \item[Ad spend] The media cost for the campaign (excludes production costs unless specified
    as fully loaded CPA).
  \item[Conversions] The specific action being measured --- a paying customer (CPA), a
    sales-qualified lead (CPL), a trial signup, a form fill, depending on the measurement
    context. The definition must be fixed before comparing CPA across campaigns.
\end{description}

\textbf{What it tracks.} The average cost of driving one target action from paid media.
CPA is the performance marketer's primary efficiency metric, sitting one level above CTR/CPC
and directly comparable to CAC when the action is a paying customer.

\textbf{Why it matters.} CPA is the link between media spend and unit economics: if
$\mathrm{CPA} < \mathrm{LTV}/3$, the channel is economically viable at the target LTV:CAC.
The ceiling for CPA is derived from LTV in \cref{ch:paid-acquisition-platforms}.

\textbf{Benchmark.} The LTV-implied ceiling: $\mathrm{CPA}_{\max} = \mathrm{LTV} / 3$ for
an LTV:CAC of 3. For the running SaaS, $\mathrm{CPA}_{\max} = \money{3150}/3 = \money{1050}$
per paying customer.

\textbf{Trap.} Defining conversions too loosely to flatter the metric. Counting form fills
(CPL at \money{20}) rather than sales-qualified leads (CPSQL at \money{200}) can make a
campaign look 10$\times$ more efficient than it actually is at the revenue level. Every CPA
figure requires a clear, stable definition of the conversion event.

\begin{example}[CPA for the running SaaS]
Google Search campaign for the month: spend \money{52000}; new paying customers: 87.
\[
  \mathrm{CPA} = \frac{\money{52000}}{87} \approx \money{598}.
\]
This sits just below the \money{600} blended CAC, confirming the Google Search channel is
contributing at roughly the average acquisition efficiency. The LTV-implied ceiling of
\money{1050} means there is significant headroom to increase bids before the channel turns
uneconomical.
\end{example}


\subsection{ROAS --- Return on Ad Spend}

\textbf{Formula.}
\[
  \mathrm{ROAS} = \frac{\text{Conversion value (revenue)}}{\text{Ad spend}}
\]

\textbf{What each term means.}
\begin{description}
  \item[Conversion value] Revenue attributed to conversions driven by the campaign, as
    reported by the ad platform. The platform's attribution model (typically last-click)
    determines what revenue is credited.
  \item[Ad spend] The media cost for the campaign in the measurement period.
\end{description}

\textbf{What it tracks.} Revenue returned per dollar of ad spend, as measured by the
platform's own attribution. ROAS is the standard efficiency metric for e-commerce and direct
response campaigns where the platform can track revenue directly.

\textbf{Why it matters.} ROAS translates media cost into a revenue multiple, allowing fast
comparison across campaigns and channels. It connects directly to profitability: a ROAS of
3 with a \qty{40}{\percent} gross margin means the campaign is breaking even on gross profit
($3 \times 0.40 = 1.20$ gross profit per dollar spent, minus the dollar of spend itself
leaves \money{0.20} of gross margin cover, which must then absorb fixed costs).

\textbf{Benchmark.} For e-commerce: 3--5$\times$ is broadly acceptable; 5--8$\times$ is
strong. The right ROAS target must be derived from gross margin: $\mathrm{ROAS}_{\min} =
1/\text{gross margin}$ for break-even on gross profit. All platform-reported ROAS figures
are subject to attribution inflation (\cref{ch:attribution-and-measurement}).

\textbf{Trap.} Platform-reported ROAS systematically over-credits the platform. Every
channel's attribution model claims credit for conversions that also appear in other channels'
reports (double-counting), and view-through credits inflate ROAS on impression-heavy
channels. Use ROAS as a campaign-management tool, not as a cross-channel efficiency
comparison; use iROAS for causal measurement.

\begin{example}[ROAS for a direct-to-consumer illustrative case]
A \textit{consumer supplement brand} (non-SaaS, for illustration) runs a Meta Advantage+
campaign: spend \money{25000}; platform-reported revenue: \money{125000}; ROAS $= 5\times$.
An incrementality test (geo holdout) reveals true incremental revenue: \money{60000};
iROAS $= 2.4\times$ --- barely above break-even at the firm's \qty{45}{\percent} gross
margin. The platform's ROAS over-credited by more than \qty{100}{\percent}.
\end{example}


\subsection{MER --- Marketing Efficiency Ratio}

\textbf{Formula.}
\[
  \mathrm{MER} = \frac{\text{Total revenue}}{\text{Total marketing spend}}
\]

\textbf{What each term means.}
\begin{description}
  \item[Total revenue] All revenue in the period --- not only attributed revenue. The
    whole-business topline.
  \item[Total marketing spend] All marketing spend across all paid channels in the same
    period.
\end{description}

\textbf{What it tracks.} The blended, whole-business efficiency of the marketing investment.
Unlike per-channel ROAS (which depends on attribution models), MER is model-free: it simply
asks how much revenue the business generates per dollar of total marketing spend.

\textbf{Why it matters.} MER sidesteps the attribution debate entirely. When channel-level
ROAS figures from Meta, Google, and TikTok sum to more than total business revenue (because
each platform over-counts), MER provides the honest sanity check. A business with \money{4}
MER knows that for every dollar of total marketing spend, the business generates \money{4} of
total revenue --- regardless of which channel ``caused'' it.

\textbf{Benchmark.} Varies widely by business model and margin structure. For SaaS:
MER of 10--20$\times$ on new-customer marketing spend is typical at scale (because LTV is
high and CAC is low relative to ACV). For e-commerce with thin margins, MER of 3--6$\times$
may be the viable range.

\textbf{Trap.} MER hides which channel actually drove revenue. A firm that turns off its
best-performing channel and shifts budget to a worse one may maintain MER in the short
run (because brand equity from the good channel sustains revenue briefly) while
systematically degrading future performance. Use MER as an aggregate health check, not a
channel allocation tool.

\begin{example}[MER for the running SaaS]
Total MRR this month: \money{340000} (all revenue, new and recurring). Total marketing
spend: \money{120000} (the blended CAC numerator).
\[
  \mathrm{MER} = \frac{\money{340000}}{\money{120000}} \approx 2.83.
\]
This MER is deliberately low because it divides total \emph{monthly} revenue by the
acquisition-only spend (most revenue is from existing customers, not new acquisition).
On a new-customer-revenue-only basis (200 new customers $\times$ \money{50}/month
$= \money{10000}$): $\mathrm{MER}_\text{acquisition} = \money{10000}/\money{120000}
= 0.08$ --- which is expected, since the business model is recurring and the payback
period is 29 months.
\end{example}


\subsection{AOV --- Average Order Value}

\textbf{Formula.}
\[
  \mathrm{AOV} = \frac{\text{Total revenue}}{\text{Number of orders}}
\]

\textbf{What each term means.}
\begin{description}
  \item[Total revenue] Revenue from all transactions in the measurement period.
  \item[Number of orders] Count of distinct purchase transactions in the same period.
\end{description}

\textbf{What it tracks.} The average size of a single transaction. AOV is primarily a
retail and e-commerce metric; in SaaS the equivalent is ARPU (average revenue per user).

\textbf{Why it matters.} AOV directly determines the maximum viable CPA: a channel that
drives buyers at an AOV of \money{30} with a \qty{40}{\percent} gross margin has at most
\money{12} of gross profit to fund acquisition. Raising AOV through upsell, cross-sell, or
bundling is one of the highest-leverage interventions in e-commerce because it improves
profitability without changing traffic volume or conversion rate.

\textbf{Benchmark.} Entirely category-dependent. AOV benchmarks only make sense within a
vertical (e.g., fashion e-commerce, B2B software seats). Trend over time matters more than
the absolute level.

\textbf{Trap.} The mean hides a long, skewed tail. A small number of large orders can pull
AOV well above the typical customer's spend, making the average useless for pricing and
bundling decisions. Always examine the median and the distribution, not only the mean.

\begin{example}[AOV for an illustrative e-commerce case]
A \textit{direct-to-consumer coffee brand} (illustrative): \num{4200} orders in the month;
revenue \money{189000}.
\[
  \mathrm{AOV} = \frac{\money{189000}}{4200} = \money{45}.
\]
The mean is \money{45} but the median order is \money{32} --- the gap driven by 200
subscription-bundle orders averaging \money{120}. Optimising advertising toward the bundle
buyers raises AOV and dramatically improves CPA viability.
\end{example}


\subsection{Conversion Rate}

\textbf{Formula.}
\[
  \mathrm{CVR} = \frac{\text{Conversions}}{\text{Sessions (or visitors)}}
\]

\textbf{What each term means.}
\begin{description}
  \item[Conversions] The specific target action: a paid purchase, a trial signup, a
    demo booked, or any other defined event. The action must match the metric's purpose.
  \item[Sessions] The number of distinct user sessions (visits) to the page or funnel step
    being measured in the same period.
\end{description}

\textbf{What it tracks.} The efficiency of a specific funnel step --- the fraction of
visitors who take the target action. CVR can be measured at each stage: session-to-trial,
trial-to-paid, demo-to-closed, and so on.

\textbf{Why it matters.} CVR links traffic volume to revenue outcomes and therefore links
CPC to CPA ($\mathrm{CPA} = \mathrm{CPC}/\mathrm{CVR}$). A \qty{1}{\percent} improvement
in conversion rate has the same effect on CPA as a \qty{1}{\percent} reduction in CPC ---
but conversion rate improvement typically costs R\&D/design time rather than media spend.
The full funnel chain is formalised as \cref{eq:funnel-chain} in \cref{ch:digital-funnel-and-crm}.

\textbf{Benchmark.} Highly context-dependent. Session-to-trial for B2B SaaS: \qty{2}{\percent}--\qty{5}{\percent}.
Trial-to-paid: \qty{15}{\percent}--\qty{30}{\percent} with active onboarding.
Site-wide e-commerce CVR: \qty{1}{\percent}--\qty{4}{\percent}.

\textbf{Trap.} Site-wide CVR is almost meaningless without segmentation by traffic source,
device, landing page, and user intent. Branded search visitors convert at 5--10$\times$ the
rate of cold display traffic; averaging them produces a number that describes neither group
accurately and misleads optimisation decisions in both directions.

\begin{example}[Conversion Rate in the running SaaS]
From the fact-sheet lock: \num{10000} monthly sessions $\to$ \num{800} trials
$\to$ \num{200} paid customers.
\[
  \mathrm{CVR}_{\text{session-to-trial}} = \frac{800}{10000} = \qty{8}{\percent}; \qquad
  \mathrm{CVR}_{\text{trial-to-paid}} = \frac{200}{800} = \qty{25}{\percent}.
\]
The session-to-trial rate is strong (above typical benchmark), suggesting good organic SEO
intent. The trial-to-paid rate is at the high end of the healthy range, indicating effective
onboarding. The constraint is top-of-funnel traffic volume, not conversion efficiency.
\end{example}


%% ============================================================
\section{Ad-Platform}

Ad-platform metrics live inside the auction mechanics of paid channels. They diagnose
the efficiency and coverage of the firm's presence within a specific platform, and they
inform bidding strategy rather than overall marketing strategy.

\subsection{eCPM --- Effective Cost Per Thousand Impressions}

\textbf{Formula.}
\[
  \mathrm{eCPM} = \frac{\text{Total cost}}{\text{Total impressions}} \times 1000
\]

\textbf{What each term means.}
\begin{description}
  \item[Total cost] All-in ad spend across whatever buying type was used --- cost-per-click,
    cost-per-action, or direct CPM --- for the campaign or period.
  \item[Total impressions] All impressions served under that spend, regardless of buy type.
\end{description}

\textbf{What it tracks.} The all-in realised cost per thousand impressions across multiple
buying types. eCPM normalises different auction mechanics (CPC, CPA, CPM buys) into a
single comparable cost-of-reach figure.

\textbf{Why it matters.} When a business runs campaigns under different buy types --- some
optimised for clicks, some for conversions, some for reach --- eCPM provides a consistent
cross-campaign comparison of how efficiently each is purchasing audience. It is also the
publisher's revenue metric: the supply side of the programmatic auction always quotes in eCPM
terms.

\textbf{Benchmark.} Use eCPM only for comparison within the same placement type and
audience segment. Premium placements (homepage takeovers, in-feed social) carry
structurally higher eCPMs than run-of-network display, making cross-placement eCPM
comparisons misleading.

\textbf{Trap.} Comparing eCPM across very different placements or audiences treats all
impressions as equivalent when they are not. A \money{1.50} eCPM on remnant display and a
\money{60} eCPM on LinkedIn's executive feed may both be appropriate for their respective
purposes; the comparison says nothing useful about relative efficiency on an outcome basis.

\begin{example}[eCPM in the running SaaS]
A mixed-buy Google campaign: \money{30000} total spend; \num{2000000} impressions delivered
across Search (\num{400000}) and Display (\num{1600000}).
\[
  \mathrm{eCPM} = \frac{\money{30000}}{2000000} \times 1000 = \money{15}.
\]
Search impressions alone: \money{20000} spend / \num{400000} impressions $\times$ 1000
$= \money{50}$ eCPM. Display alone: \money{10000} / \num{1600000} $\times$ 1000
$= \money{6.25}$ eCPM. The blended \money{15} masks the structural difference --- compare
Search eCPM only against other Search campaigns, and Display only against Display.
\end{example}


\subsection{Impression Share}

\textbf{Formula.}
\[
  \text{Impression Share} = \frac{\text{Impressions won}}{\text{Impressions eligible}}
\]

\textbf{What each term means.}
\begin{description}
  \item[Impressions won] The number of auctions in which the ad was actually shown to a user.
  \item[Impressions eligible] The estimated number of auctions the ad was eligible to enter
    based on its targeting settings, budget, and quality signals.
\end{description}

\textbf{What it tracks.} What fraction of available eligible impressions the campaign is
actually capturing. Impression share below \qty{100}{\percent} means the campaign is missing
auctions it is qualified to win --- either because the budget ran out (budget-limited IS) or
because the ad's quality or bid rank was too low (rank-limited IS).

\textbf{Why it matters.} Impression share diagnoses the bottleneck in paid coverage. If IS
is low due to budget, adding budget will directly increase coverage. If IS is low due to rank
(ad quality or bid), increasing budget does nothing; Quality Score improvement or bid
increases are required. Splitting lost-IS into its two components (budget vs rank) before
changing strategy prevents the common error of adding budget to a rank-limited campaign.

\textbf{Benchmark.} Healthy target depends on competitive intensity. Chasing
\qty{100}{\percent} IS on broad keywords is a budget trap; \qty{80}{\percent}+ IS on
the highest-intent, highest-converting terms is a more defensible goal.

\textbf{Trap.} Targeting \qty{100}{\percent} impression share at any price. At the margin,
winning every possible impression requires outbidding all competitors on all terms, including
low-intent terms that convert poorly. The incremental impressions at the margin cost more
and convert less than the average --- classic diminishing returns to IS maximisation.

\begin{example}[Impression Share in the running SaaS]
Google Search campaign targeting ``project management software'' and five related terms.
Impressions won: \num{32000}; estimated eligible impressions: \num{50000}.
\[
  \text{IS} = \frac{32000}{50000} = \qty{64}{\percent}.
\]
Lost IS breakdown: \qty{20}{\percent} lost to budget, \qty{16}{\percent} lost to rank.
Strategy: increase budget to capture the \qty{20}{\percent} budget-lost IS first (immediate,
controllable) and run ad-copy tests to improve Quality Score for the rank-lost share (slower,
but reduces average CPC).
\end{example}


\subsection{ACOS --- Advertising Cost of Sales}

\textbf{Formula.}
\[
  \mathrm{ACOS} = \frac{\text{Ad spend}}{\text{Attributed sales revenue}}
\]

\textbf{What each term means.}
\begin{description}
  \item[Ad spend] The advertising cost for the product or campaign measured.
  \item[Attributed sales revenue] Sales revenue attributed to the ad, typically within a
    defined attribution window. ACOS is the primary metric on Amazon's ad platform; the
    attribution window is platform-defined.
\end{description}

\textbf{What it tracks.} Ad cost expressed as a share of attributed sales --- the inverse
of ROAS. ACOS of \qty{25}{\percent} means \money{0.25} was spent on ads for every
\money{1.00} of sales.

\textbf{Why it matters.} ACOS is the native efficiency metric for Amazon Sponsored Products
and similar retail-media platforms. It is used to determine whether a product's ad spend is
within the breakeven margin. Because it is the inverse of ROAS, the interpretation is
``smaller is better'' --- unlike ROAS where higher is better.

\textbf{Benchmark.} Target ACOS $=$ product gross margin. If the product margin is
\qty{35}{\percent}, a target ACOS of \qty{35}{\percent} means ad spend equals gross profit
--- the product breaks even at the contribution level. ACOS below the margin target
generates gross profit from advertising; above it destroys margin.

\textbf{Trap.} Setting ACOS targets without anchoring to product margin. An ACOS of
\qty{20}{\percent} on a product with \qty{15}{\percent} margin means every advertised sale
destroys money; the same \qty{20}{\percent} ACOS on a \qty{50}{\percent} margin product is
highly profitable. The target must be derived from the specific product's economics.

\begin{example}[ACOS for an illustrative Amazon seller]
A \textit{consumer hardware brand} (illustrative, not SaaS): product sells at \money{89};
COGS \money{32}; gross margin \qty{64}{\percent}. Sponsored Products spend: \money{4450};
attributed sales: \money{22250}.
\[
  \mathrm{ACOS} = \frac{\money{4450}}{\money{22250}} = \qty{20}{\percent}.
\]
Since \qty{20}{\percent} $<$ \qty{64}{\percent} margin, the campaign generates
$(\qty{64}{\percent} - \qty{20}{\percent}) = \qty{44}{\percent}$ gross margin on advertised
sales --- profitable advertising. The breakeven ACOS is \qty{64}{\percent}; the business has
significant headroom before the campaign destroys value.
\end{example}


\subsection{iROAS --- Incremental Return on Ad Spend}

\textbf{Formula.}
\[
  \mathrm{iROAS} = \frac{\text{Incremental revenue caused by the ad}}{\text{Ad spend}}
\]

\textbf{What each term means.}
\begin{description}
  \item[Incremental revenue] Revenue that \emph{would not have occurred} without the
    advertising --- measured via a holdout experiment (geo test, user holdout, or switchback)
    that isolates causation from correlation.
  \item[Ad spend] The cost of the advertising in the test cell.
\end{description}

\textbf{What it tracks.} The causal return on advertising spend --- how much additional
revenue the ad genuinely caused, versus revenue that would have happened organically or was
attributed by the platform's last-click model regardless of the ad's role.

\textbf{Why it matters.} iROAS is the gold standard for advertising measurement because it
answers the only question that matters for budget allocation: if we spend \money{1} more on
this channel, does the business get more than \money{1} back? Platform-reported ROAS cannot
answer this question because it attributes organic and assisted conversions to paid media. The
measurement methodology is fully developed in \cref{ch:attribution-and-measurement}.

\textbf{Benchmark.} iROAS must exceed 1 for the spend to be cash-flow positive. It must
exceed $1/\text{gross margin}$ for the spend to be gross-profit positive. For the running
SaaS with gross margin approximately \qty{79}{\percent}, $\mathrm{iROAS}_{\min} = 1/0.79
\approx 1.27$.

\textbf{Trap.} iROAS requires a holdout or geo test --- it cannot be read from platform
dashboards. Without an experiment, iROAS is not observable. Claims of iROAS derived from
attribution models (even sophisticated MMM models) are model outputs, not causal
measurements; they carry model uncertainty that must be disclosed alongside the figure.

\begin{example}[iROAS for the running SaaS]
A geo holdout test: test markets receive \money{40000} in incremental brand spend over a
six-week period; control markets receive normal spend. Test-market revenue: \money{1800000};
counterfactual control-market revenue (scaled): \money{1620000}. Incremental revenue:
\money{180000}.
\[
  \mathrm{iROAS} = \frac{\money{180000}}{\money{40000}} = 4.5.
\]
Platform-reported ROAS for the same campaign: $7\times$ --- a \qty{56}{\percent}
overstatement. The iROAS of 4.5 still clears the gross-margin floor of 1.27, confirming the
brand spend creates value; but the budget decision would have been distorted by the
platform-reported figure.
\end{example}


%% ============================================================
\section{Product \& Engagement}

Product metrics measure whether users are finding lasting value in the product, not merely
visiting it. An acquisition engine that fills a leaky product is not a growth engine ---
it is a churn factory. These metrics diagnose the product's hold on user behaviour and
the durability of the customer relationship.

\subsection{Stickiness --- DAU/MAU Ratio}

\textbf{Formula.}
\[
  \text{Stickiness} = \frac{\mathrm{DAU}}{\mathrm{MAU}}
\]

\textbf{What each term means.}
\begin{description}
  \item[DAU] Daily Active Users: the count of unique users who performed at least one
    meaningful action in the product on a given day. ``Meaningful'' must be defined (a
    login is often too weak a signal; a core-feature interaction is more defensible).
  \item[MAU] Monthly Active Users: the count of unique users who performed at least one
    meaningful action in the product at any point in the calendar month.
\end{description}

\textbf{What it tracks.} The fraction of monthly active users who return on a daily basis
--- a direct proxy for how habitual and embedded the product has become in users' daily
workflows.

\textbf{Why it matters.} Stickiness predicts retention. A product with \qty{50}{\percent}
stickiness is genuinely woven into daily routine; at \qty{10}{\percent}, users visit
occasionally and are one bad experience from churning. For B2B SaaS, stickiness also drives
expansion revenue: teams that use a product daily are far more likely to add seats and
upgrade than occasional users.

\textbf{Benchmark.}
\begin{itemize}
  \item $>\qty{50}{\percent}$: elite; characteristic of daily-utility products (Slack,
    Gmail, Notion).
  \item $>\qty{20}{\percent}$: decent for B2B tools not designed for daily use (quarterly
    planning, annual reporting tools naturally sit lower).
  \item $<\qty{10}{\percent}$: low engagement; typically predicts high churn and low
    expansion.
\end{itemize}

\textbf{Trap.} Bot accounts and duplicate user accounts inflate DAU and MAU independently,
but the ratio can still look healthy even as real user engagement falls. Always validate
active-user counts against login and action deduplication before reporting stickiness.

\begin{example}[Stickiness in the running SaaS]
In the measurement month: MAU $=$ \num{5500} (out of \num{6700} paying customers, some
teams share logins or use the product infrequently). Average daily active users across
the month: \num{1650}.
\[
  \text{Stickiness} = \frac{1650}{5500} = \qty{30}{\percent}.
\]
This is above the B2B decent threshold. The \num{1200} paying customers not in MAU are
flagged for a proactive customer-success outreach --- low-usage accounts are the leading
indicator of churn, not trailing.
\end{example}


\subsection{D30 Retention --- Day-30 Retention Rate}

\textbf{Formula.}
\[
  \text{D30 Retention} = \frac{\text{Users active at day 30}}{\text{Cohort size at day 0}}
\]

\textbf{What each term means.}
\begin{description}
  \item[Users active at day 30] Users from the acquisition cohort who performed at least one
    meaningful action in the product on day 30 (or within a defined window around day 30,
    e.g., days 27--33).
  \item[Cohort size at day 0] The number of users who joined (or activated) in the cohort.
    A ``cohort'' is a group defined by a shared starting event --- typically the signup week
    or the first-login week.
\end{description}

\textbf{What it tracks.} The fraction of a new-user cohort that is still engaging with the
product 30 days after joining. D30 retention is the first durable signal of whether the
product has created lasting value for new users --- early-day retention (D1, D7) is often
inflated by novelty.

\textbf{Why it matters.} D30 retention is the strongest early predictor of long-run customer
LTV. A cohort that retains at \qty{40}{\percent} on day 30 will have a very different
long-run retention curve than one that retains at \qty{15}{\percent}. It also diagnoses
the onboarding funnel: a D30 drop-off concentrated in the first two weeks points to an
onboarding problem; a drop-off in weeks three and four points to a value-realisation problem
(users got on the product but did not find its core value). Full retention curve analysis
is in \cref{ch:digital-funnel-and-crm}.

\textbf{Benchmark.} Highly product-category-dependent. Consumer mobile: \qty{20}{\percent}
D30 retention is a reasonable floor. B2B SaaS with active onboarding: \qty{60}{\percent}+
D30 retention is achievable and important for the unit economics to work.

\textbf{Trap.} Averaging D30 retention across acquisition sources masks the quality
difference between channels. Organic SEO cohorts may retain at \qty{70}{\percent} while
paid social cohorts from a broad-audience campaign retain at \qty{25}{\percent}; the blended
figure of \qty{45}{\percent} describes neither cohort and sends the acquisition team a
misleading signal about paid channel quality.

\begin{example}[D30 Retention in the running SaaS]
A February signup cohort of \num{800} trial users (from \num{10000} sessions $\to$
\num{800} trials). Active on day 30: \num{520}.
\[
  \text{D30 Retention} = \frac{520}{800} = \qty{65}{\percent}.
\]
Within this cohort, Google-organic signups (600): day-30 retention \qty{74}{\percent}.
Paid-LinkedIn signups (200): day-30 retention \qty{41}{\percent}. The blended figure of
\qty{65}{\percent} conceals a meaningful quality difference that should inform how the
\money{120000} monthly acquisition budget is allocated.
\end{example}


\subsection{NPS --- Net Promoter Score}

\textbf{Formula.}
\[
  \mathrm{NPS} = \%\,\text{Promoters} - \%\,\text{Detractors}
\]
where Promoters are respondents who score 9--10 on the ``recommend'' question, and Detractors
score 0--6.

\textbf{What each term means.}
\begin{description}
  \item[\% Promoters] Share of survey respondents giving a score of 9 or 10 out of 10 on
    the question: ``How likely are you to recommend [product] to a colleague?''
  \item[\% Detractors] Share of respondents giving a score of 0--6 (inclusive). Scores of
    7--8 are ``Passives'' and are excluded from the calculation.
\end{description}

\textbf{What it tracks.} A coarse, compressed proxy for customer satisfaction and word-of-mouth
propensity. NPS ranges from $-100$ (all Detractors) to $+100$ (all Promoters).

\textbf{Why it matters.} NPS is widely used as a top-line customer sentiment indicator and
for benchmarking against industry peers. It is predictive of churn at the cohort level ---
Detractor cohorts churn at materially higher rates than Promoter cohorts --- and it
captures the word-of-mouth component of the growth model.

\textbf{Benchmark.} Industry-dependent. SaaS: above 30 is positive, above 50 is strong.
Below 0 is a warning sign requiring immediate qualitative investigation.

\textbf{Trap.} Reading an 11-point survey as a precise number. A shift in NPS from 42 to 47
in a sample of 300 respondents is statistically indistinguishable from noise --- the
confidence interval typically spans 10--15 points. Trend direction over quarters matters
more than point-in-time precision. Also: survey response bias systematically over-represents
highly satisfied and highly dissatisfied customers, missing the ambivalent middle.

\begin{example}[NPS in the running SaaS]
Quarterly NPS survey sent to \num{6700} customers; \num{1340} responses (\qty{20}{\percent}
response rate). Promoters (9--10): 670 (\qty{50}{\percent}). Passives (7--8): 469
(\qty{35}{\percent}). Detractors (0--6): 201 (\qty{15}{\percent}).
\[
  \mathrm{NPS} = 50 - 15 = 35.
\]
The 201 Detractor accounts are cross-referenced against usage data: \qty{78}{\percent} of
them have stickiness below \qty{10}{\percent}, confirming low engagement as the primary
dissatisfaction driver and prioritising the customer-success intervention.
\end{example}


\subsection{Activation Rate}

\textbf{Formula.}
\[
  \text{Activation Rate} = \frac{\text{Users who reach the ``aha moment''}}{\text{Total signups}}
\]

\textbf{What each term means.}
\begin{description}
  \item[``Aha moment''] The specific in-product action or milestone that correlates most
    strongly with long-run retention. For a project-management SaaS, this might be
    ``created first project with 3+ collaborators.'' For CRM: ``logged first deal and set
    a follow-up.'' The aha moment is defined empirically from cohort retention analysis, not
    assumed.
  \item[Total signups] All users who registered during the measurement period, including
    trials and free accounts.
\end{description}

\textbf{What it tracks.} The fraction of new users who successfully reach the product's core
value moment during onboarding. Activation rate is the hinge between acquisition and
retention: users who activate retain at 3--5$\times$ the rate of those who do not.

\textbf{Why it matters.} Improving activation rate multiplies the value of every acquisition
dollar: if 30 of every 100 trial signups activate, doubling activation to 60 effectively
doubles the yield of acquisition without spending an extra dollar on media. It is the
highest-leverage onboarding intervention available. The activation funnel connects directly
to \cref{ch:digital-funnel-and-crm}.

\textbf{Benchmark.} Strong activation rates vary by product complexity. Simple consumer
tools: \qty{60}{\percent}+ is achievable. Complex B2B platforms: \qty{30}{\percent}--\qty{50}{\percent}
is healthy; below \qty{20}{\percent} typically signals an onboarding problem severe enough
to override acquisition investment.

\textbf{Trap.} Defining ``activated'' too loosely inflates the metric without predicting
retention. If activation is defined as ``logged in at least once,'' nearly all signups
``activate'' but the metric predicts nothing. The correct definition is the event that
maximally differentiates the long-run retention curves of activated vs non-activated users.

\begin{example}[Activation Rate in the running SaaS]
The product team defines activation as: ``created a workspace, added at least two
integrations, and invited one collaborator within the first 7 days.'' Of \num{800} trial
signups this month, \num{296} reached this milestone.
\[
  \text{Activation Rate} = \frac{296}{800} = \qty{37}{\percent}.
\]
Retention analysis shows activated users retain at \qty{82}{\percent} on day 30; non-activated
users at \qty{19}{\percent}. The 37\% activation rate means 504 trial users are on the
low-retention path. A targeted onboarding email sequence with a \qty{30}{\percent} lift in
activation would add approximately 151 activated users per month, raising expected paid
conversions from 200 to 238 at no additional acquisition spend --- equivalent to a
\qty{19}{\percent} increase in effective marketing ROI.
\end{example}


%% ============================================================
\section{Vanity Metrics --- What They Hide}

Vanity metrics move with effort, not with outcomes. Each one in this section actively
misleads decision-making rather than merely being imprecise: it can increase while the
business deteriorates, and optimising for it accelerates the deterioration. Each metric
below carries an \textbf{admonition} in place of a worked example, plus a pointer to the
outcome metric that should replace it.

\subsection{Raw Impressions}

\textbf{Formula.}
\[
  \text{Raw Impressions} = \text{total ad views (served impressions)}
\]

\textbf{What each term means.}
\begin{description}
  \item[Served impressions] The number of times an ad was technically delivered to a device.
    This includes impressions never actually seen by a human (below-the-fold, auto-play
    with no viewability verification, invalid traffic).
\end{description}

\textbf{What it tracks.} Spend, not reach. Raw impression volume scales mechanically with
budget; it does not indicate that a real human saw the ad, that the ad was relevant, or that
it produced any effect.

\textbf{Why it misleads.} A campaign can multiply impression count by increasing budget,
loosening targeting, or buying cheap remnant inventory --- none of which improves business
outcomes. Use \emph{unique reach with viewability} (the number of distinct humans who saw a
verified-viewable impression) and \emph{frequency} (impressions per unique user) instead.

\begin{admonition}{Warning}
Raw impressions can increase by \qty{300}{\percent} while effective reach remains flat ---
simply by purchasing the same audience three times rather than reaching three times as many
people. A campaign reporting ``500 million impressions'' could represent 50 million people
each seeing the ad 10 times (excessive frequency, likely with diminishing returns by
impression 4) or 500 million different people each seeing it once (broad, light-touch
reach). The impression count is identical; the business implication is opposite.
Use reach and frequency as the paired reporting unit, never impressions alone.
\end{admonition}


\subsection{Email Open Rate}

\textbf{Formula.}
\[
  \text{Open Rate} = \frac{\text{Opens}}{\text{Emails delivered}}
\]

\textbf{What each term means.}
\begin{description}
  \item[Opens] Traditionally tracked by a 1-pixel tracking image loaded when the email is
    displayed. Since Apple's Mail Privacy Protection (MPP) launched in September 2021,
    Apple Mail pre-fetches this pixel on behalf of the user whether or not they opened the
    email.
  \item[Emails delivered] Sent minus hard bounces.
\end{description}

\textbf{What it tracks.} Almost nothing since Apple MPP. For lists with significant Apple
Mail usage (\qty{50}{\percent}+), reported open rates are largely a function of how many
subscribers use Apple Mail, not of whether they actually read the email.

\textbf{Why it misleads.} Open rates reported in most ESPs have been systematically inflated
since September 2021 by Apple's privacy proxies. A list that shows \qty{48}{\percent} open
rate may have a true human-open rate of \qty{20}{\percent}. Optimising subject lines for
open rate using this data optimises for what the Apple proxy fetches, not for human
behaviour. Use \emph{click rate} (clicks per delivered email) and \emph{reply rate} as the
primary engagement metrics; for direct-response emails, use \emph{revenue per email sent}.

\begin{admonition}{Warning}
A B2B SaaS measured email performance by open rate and celebrated a rise from \qty{32}{\percent}
to \qty{47}{\percent} after a subject-line test. Click rates were flat. Revenue per email
sent was flat. The open-rate lift was entirely attributable to a shift in list composition
toward Apple Mail users --- the proxy inflated the denominator at a different rate. Decisions
made on this data reallocated budget toward email and away from channels with measurable
click-and-conversion data, harming overall marketing ROI. Report open rates only as a
historical artefact; never make resource-allocation decisions from them.
\end{admonition}


\subsection{Last-Click ROAS}

\textbf{Formula.}
\[
  \text{Last-Click ROAS} = \frac{\text{Revenue attributed to last click}}{\text{Ad spend}}
\]

\textbf{What each term means.}
\begin{description}
  \item[Revenue attributed to last click] Revenue from conversions where the reported
    channel received the final click before purchase. All upstream channels (awareness,
    consideration, retargeting assists) receive zero credit.
  \item[Ad spend] The cost of the campaign being measured.
\end{description}

\textbf{What it tracks.} Who closed the sale, not who caused the sale. Last-click ROAS
identifies the channel that happened to be in contact at the moment of conversion, which in
a multi-touch journey is typically a branded-search or retargeting campaign that intercepted
demand generated by an earlier touchpoint.

\textbf{Why it misleads.} Last-click attribution systematically over-credits the final
touchpoint (branded search, direct, retargeting) and under-credits the channels that created
initial awareness and intent (social prospecting, content, PR). Budget optimised on
last-click ROAS shifts spend toward demand-\emph{capturing} channels and defunds
demand-\emph{creating} channels, eventually starving the top of the funnel and causing
last-click ROAS itself to decay as the underlying awareness base erodes. Use
\emph{iROAS from incrementality experiments} as the causal alternative
(\cref{ch:attribution-and-measurement}).

\begin{admonition}{Warning}
A well-known e-commerce brand optimised exclusively on last-click ROAS for 18 months.
Branded-search ROAS rose from $6\times$ to $11\times$ (it was capturing demand created by
a brand equity built over years). Simultaneously, top-of-funnel spend was cut from
\qty{35}{\percent} to \qty{12}{\percent} of budget. New-customer acquisition volume fell
\qty{40}{\percent} over the period; the existing customer base aged out and revenue declined
two years later. The last-click ROAS number was the highest it had ever been on the day the
brand's organic search traffic began its structural decline. iROAS tests, run too late,
showed the top-of-funnel channels had an iROAS of 3.8 --- highly profitable and starved to
zero.
\end{admonition}


\subsection{Unqualified Leads}

\textbf{Formula.}
\[
  \text{Unqualified Leads} = \text{form fills or contact requests}
\]

\textbf{What each term means.}
\begin{description}
  \item[Form fills] Any submission of a web form, gated-content download, demo-request
    click, or contact request, regardless of buyer intent, company fit, or purchase
    authority.
\end{description}

\textbf{What it tracks.} Interest, not demand. An unqualified lead has expressed some form
of interest but has not been evaluated against the criteria that make a prospect a viable
buyer (company size, budget, authority, need, timing --- BANT or equivalent framework).

\textbf{Why it misleads.} Unqualified lead volume scales with top-of-funnel spend and
gated-content volume, both of which can be increased mechanically with no improvement in
pipeline quality. A content team that publishes a viral ``Free ROI Calculator'' download
may generate 5,000 form fills in a month with zero intent to buy. Reporting lead volume to
leadership creates pressure to generate more leads, not better ones, resulting in sales teams
buried in low-quality outreach. Feed \emph{qualified pipeline value} (sales-qualified
opportunities with attached revenue) back to the marketing team instead, and attribute
revenue against spend rather than leads against spend. See
\cref{ch:paid-acquisition-platforms} for the full CPL-to-CPA conversion mechanics.

\begin{admonition}{Warning}
A B2B SaaS marketing team was measured on monthly lead volume. To hit 2,000 leads per month,
the team purchased a third-party email list and promoted a generic whitepaper. Leads hit
2,500 --- a record. Sales accepted leads (SALs): 48. Closed won: 3. Cost per closed customer
from this campaign: over \money{15000}, versus the blended CAC of \money{600}. The lead
number rose; the business was harmed. Measure \emph{qualified pipeline generated} and
\emph{revenue influenced} instead --- metrics the sales and marketing teams must agree on
before the campaign launches.
\end{admonition}


\subsection{Follower Count}

\textbf{Formula.}
\[
  \text{Follower Count} = \text{total followers on a social platform}
\]

\textbf{What each term means.}
\begin{description}
  \item[Followers] User accounts that have clicked ``Follow'' on the brand's social profile.
    On most platforms, the majority of followers see less than \qty{5}{\percent} of organic
    posts due to algorithmic distribution limits.
\end{description}

\textbf{What it tracks.} Audience size on a platform, not audience value. Follower count
measures a historical accumulation of ``Follow'' clicks --- it does not measure engagement,
reach of current content, or any revenue linkage.

\textbf{Why it misleads.} Follower count can be grown through follow/unfollow tactics,
paid follower acquisition, viral off-topic content, and giveaway campaigns --- all of which
add followers with no commercial relationship to the brand. A brand with one million
followers reaching \qty{0.5}{\percent} organically (5,000 people) may be less effective than
a competitor with 50,000 highly engaged followers reaching \qty{20}{\percent} (10,000
people). Track \emph{engaged reach per post}, \emph{conversion rate from social traffic},
and \emph{revenue attributable to social content} instead.

\begin{admonition}{Warning}
A B2B software company grew its LinkedIn follower count from 8,000 to 42,000 over six months
through a viral ``industry meme'' content series. Website traffic from LinkedIn: flat.
Demo requests from LinkedIn: unchanged. The followers acquired were largely individual
contributors in adjacent industries who would never purchase the product. Content strategy
was subsequently restructured around thought-leadership posts targeting the actual buyer
persona (VP Operations, Chief of Staff), reducing follower growth rate but increasing
demo-request attribution from LinkedIn by \qty{220}{\percent}.
\end{admonition}


\subsection{App Downloads}

\textbf{Formula.}
\[
  \text{App Downloads} = \text{total installs from app store}
\]

\textbf{What each term means.}
\begin{description}
  \item[Downloads / Installs] The count of times the application was installed on a device.
    This includes installs that are immediately abandoned, installs on devices that are never
    opened, and re-installs by the same user after deleting the app.
\end{description}

\textbf{What it tracks.} Trial, not adoption. A download is the beginning of the customer
journey, not a point on the retention curve.

\textbf{Why it misleads.} App store download counts can be inflated by store-optimisation
tricks, broad-audience install campaigns, and incentivised installs (``download for a
prize''). A consumer app with 10 million downloads and \qty{5}{\percent} day-30 retention
has 500,000 users; the same app with 1 million downloads and \qty{40}{\percent} day-30
retention has 400,000 users. The first looks 10$\times$ larger; the businesses are
comparable in active users. Report \emph{activation rate}, \emph{day-30 retention}, and
\emph{monthly active users} alongside download counts; never use downloads as a success
metric in isolation.

\begin{admonition}{Warning}
A mobile gaming startup celebrated 5 million downloads in its first three months, attracting
a seed extension at a \money{15000000} valuation anchored to download velocity. Day-30
retention was \qty{3}{\percent}; MAU at month four: \num{48000}. By month eight, DAU had
fallen below \num{10000} and new-install velocity had stalled as the broad-audience install
campaign exhausted the addressable market. The vanity metric drove a funding decision on
false premises; the outcome metric (D30 retention $\times$ cohort size) was available from
day one and predicted the trajectory accurately.
\end{admonition}
